\documentclass[../main]{subfiles}
\begin{document}
\chapter{Transmisión de calor}
\img{images/01/tcalor.jpg}{Las llamas de color azul,son las que se
presentan temperaturas más elevadas,superior a los $1300^\circ C$.}

\info{
Contenido}{
  Transferencia de calor, escala de temperaturas, Calorimetría
}

\actividad{Contestar las siguientes preguntas}
\begin{enumerate}
  \item Defina cuerpo, materia, energía, trabajo. Propiedades de la
    materia, estados de la materia, defina calor y temperatura,
    cambios de estado.
  \item Explique las escalas de temperatura y de la definición de calorimetría.
  \item ¿Qué es transferencia de calor? Tipos de transferencias de
    calor. Dibujar.
  \item Explicar ejemplos de cada uno de los casos de transferencia
    de calor y dibujar.
  \item Defina calor específico, capacidad calorífica, equivalente
    mecánico del calor, conductividad térmica y revestimientos.
\end{enumerate}

\actividad{Resolver los siguientes problemas}
\begin{enumerate}
  \item Conversión de temperatura:
    \begin{enumerate}
      \item Expresar en grados Fahrenheit la temperatura de $t=40^\circ C$
      \item Expresar en Celsius la temperatura de $t=122^\circ F$.
      \item Expresa en $^\circ K$ la temperatura del ejercicio anterior.
      \item Expresa en $^\circ F$ la temperatura de $t=746^\circ K$.
    \end{enumerate}

    \info{Calorimetría}{Se refiere a la medición de la cantidad de
      calor, sus unidades son el calor normal, la caloría o gramo caloría,

      $kcal$: la kilo caloría  o simplemente caloría.

      $BTU$: unidad que se utiliza para elevar la temperatura de
      $1^\circ F$ en una libra de agua.

    }

  \item Determinar cuantas calorías representan $45 BTU$
  \item Determinar cuantos $Chu$ representan $458 BTU$

  \item Calcular la cantidad de calor que suministra un bloque de
    mármol de 5T cuya temperatura desciende de 50C a 10C.

    \info{Calor especifico}{
      \begin{equation}
        \Delta Q= c_m G \Delta t
      \end{equation}

      Siendo

      $\Delta Q$: Calor transmitido. [Kcal]

      $c_m$: Calor específico medio $\left[\dfrac{Kcal}{Kg ^\circ C }\right]$

      $G$: masa o peso $[Kgf o Kg]$

      $\Delta t$:Diferencia de temperatura [$^\circ C$]
    }
  \item Calcular la cantidad de calor que es necesaria suministrar a
    $V=100L$ de agua para que su temperatura se eleve de 15C a 65C.
    Determinar el consumo eléctrico.

    \info{Calor transferido}{
      \begin{equation}
        Q=\dfrac{S\cdot Ct (t_1-t_2)\cdot \Delta t}{e}
      \end{equation}

      Siendo:

      $Q$: Cantidad de calor. [Kcal]

      $Ct$: Coeficiente térmico.

      $S$: Superficie. [$m^2$]

      $t_1-t_2$: Variación de temperatura. [$^\circ C$]

      $\Delta t$: tiempo [s]

      $e$: Espesor [metros]
    }

  \item Calcular la cantidad de calor que se transmite por
    conductividad a través de una losa de hormigón armado con una
    superficie de $S=16m^2$ y un espesor $e=10cm$ cuando la
    temperatura exterior es de $t_e=45^\circ C$ y la interior es
    $t_i=15^\circ C$
  \item Calcular la cantidad de calor que se transmite por
    conductividad a través del cilindro de un motor a explosión de
    acero fundido con diámetro interior de $d_i=150mm$ y diámetro
    exterior $d_e=180mm$ y $l=250mm$ de longitud, si la temperatura
    interna des de $t_i=300^\circ C$ y la externa $t_e= 90^\circ C$
  \item Calcula $Q$ que en un día pasa a través de una pared vertical
    con superficie $S=20m^2$, constituida por dos muros de
    mampostería de 20cm separados por una capa de corcho de 1cm,
    cuando la temperatura exterior es de $t_e=40^\circ C$ y la
    interior $t_i=-10^\circ C$

    \info{Paredes múltiples}{
      Cuando tenemos paredes múltiples ya sean en serie o en
      paralelo, se puede buscar una analogía con la ley de Ohm, en
      este caso la resistencia térmica. El calor transferido es
      proporcional a la diferencia de temperatura e inversamente
      proporcional a la resistencia térmica.
      \begin{equation}
        R_t=\dfrac{e}{S\cdot Ct}
      \end{equation}
      \begin{equation}
        Q=\dfrac{ \Delta t}{R_t}
      \end{equation}

      \begin{equation}
        Q=\dfrac{ \Delta t}{\dfrac{e}{S\cdot Ct}}
      \end{equation}
      \begin{equation}
        R_t=R_1+R_2+\cdots+R_n=\dfrac{e_1}{S_1\cdot
        Ct_1}+\dfrac{e_2}{S_2\cdot Ct_2}+\cdots+\dfrac{e_n}{S_n\cdot Ct_n}
      \end{equation}
      \begin{equation}
        Q=\dfrac{ \Delta t}{\dfrac{e_1}{S_1\cdot
        Ct_1}+\dfrac{e_2}{S_2\cdot Ct_2}+\cdots+\dfrac{e_n}{S_n\cdot Ct_n}}
      \end{equation}
    }
  \item Calcular Q en una pared con una superficie de $S=25m^2$
    separados por un aislamiento de hormigón ($CT_{homrigon}=1.3$)
    armado de 20cm de espesor recubiertos con ladrillos refractarios
    de 20cm ($CT_{ladrillos}=0.6$) en ambos lados y revocado en ambos
    lados de 5cm de espesor cada uno cuando la temperatura va de
    $35^\circ C$ a $10^\circ C$. Representar la Gráfica.
  \item Calcular $Q$ que en 24 horas pasa a través de una losa de
    hormigón armado de $e=25cm$ de espesor y $S=90m^2$ de superficie
    con una mampostería de ladrillo en la parte exterior de $7cm$ y
    yeso en la parte interior de $1cm$, La temperatura varía entre
    $t_i=45^\circ C$ y $t_f=0 ^\circ C$.
\end{enumerate}
\end{document}
